\documentclass[letterpaper,10pt,onecolumn]{article}
\usepackage[spanish]{babel}
\usepackage[utf8]{inputenc}
\usepackage{amsfonts}
\usepackage{amsthm}
\usepackage{amsmath}
\usepackage{mathrsfs}
%\usepackage{empheq}
\usepackage{enumitem}
\usepackage[pdftex]{color,graphicx}
\usepackage{hyperref}
\usepackage{listings}
\usepackage{calligra}
%\usepackage{algpseudocode} 
\DeclareMathAlphabet{\mathcalligra}{T1}{calligra}{m}{n}
\DeclareFontShape{T1}{calligra}{m}{n}{<->s*[2.2]callig15}{}
\newcommand{\scripty}[1]{\ensuremath{\mathcalligra{#1}}}
\lstloadlanguages{[5.2]Mathematica}
\setlength{\oddsidemargin}{0cm}
\setlength{\textwidth}{490pt}
\setlength{\topmargin}{-40pt}
\addtolength{\hoffset}{-0.3cm}
\addtolength{\textheight}{4cm}

%%%%%%%%%%%%%%%%%%%%%%%%%

\usepackage{marvosym} %use this pack
\usepackage{xcolor} %for color
%\usepackage{dtklogos}

%%%%%%%%%%%%%%%%%%%%%%%%%

\begin{document}
\begin{center}

\includegraphics[width=490pt]{header.png}\\[0.5cm]

\textsc{\huge Herramientas Computacionales}\\[0.1cm]

\Large \textsc{Contenido Programático}\\[0.7cm]

\end{center}

\large \noindent\textsc{Nombre del curso:} Herramientas Computacionales
	 
\noindent\textsc{Código del curso:} FISI 2026

\noindent\textsc{Unidad académica:} Departamento de Física

\noindent\textsc{Prerrequisitos:} Algorítmica y Programación Orientada por Objetos 1 (ISIS 1204)

\noindent\rule{\textwidth}{1pt}\\[-0.1cm]

\newcounter{mysection}
\addtocounter{mysection}{1}

\noindent\textbf{\large \Roman{mysection} \quad Introducción}\\[-0.2cm]

\noindent\normalsize Los computadores nos ayudan a organizar,
comunicar y procesar información.
Hoy en día son esenciales en los mundos de la ciencia, la academia y
la técnica. 
Este curso enseña algunas herramientas computacionales básicas para
hacer de los computadores herramientas útiles, poderosas y
versátiles para estudiantes de ciencias naturales. 
El curso desarrolla habilidades de programación en un lenguaje de alto
nivel, por ejemplo Python o Matlab; enseña algunos métodos de análisis
numérico; y exhibe algunas herramientas útiles en el análisis de 
datos. \\[0.1cm] 

\stepcounter{mysection}
\noindent\textbf{\large \Roman{mysection} \quad Objetivos}\\[-0.2cm]

\noindent\normalsize Los objetivos del curso son:

\begin{itemize}
	\item Ofrecer herramientas computacionales básicas útiles en la investigación y la vida académica.\\[-0.6cm]
	\item Introducir rutinas sencillas de análisis numérico.\\[-0.6cm]
	\item Desarrollar habilidades de programación.\\[-0.2cm]
\end{itemize}

\stepcounter{mysection}
\noindent\textbf{\large \Roman{mysection} \quad Competencias}\\[-0.2cm]

\noindent\normalsize Al finalizar el curso, se espera que el estudiante esté en capacidad de:

\begin{itemize}
	\item Utilizar computadores con sistema operativo tipo UNIX.\\[-0.6cm]
	\item Implementar en un lenguaje de programación de alto nivel la solución de problemas computacionales sencillos.\\[-0.6cm]
	\item Manipular, analizar y visualizar datos usando un lenguaje de programación de alto nivel.\\[-0.2cm]
\end{itemize}

\stepcounter{mysection}
\noindent\textbf{\large \Roman{mysection} \quad Contenido}\\[-0.2cm]

El siguiente enlace contiene los videos referenciados en el programa:

\url{https://www.youtube.com/playlist?list=PLHQtzvthdVM_MGC9dPFKe4hPAwBd_7RJ3}


\noindent\normalsize \textbf{\textsc{Semana 1}} Festivo.
\\[-0.3cm] 

\noindent\normalsize \textbf{\textsc{Semana 2}} Introducción.
Presentación. Uso de binder.
\\[-0.3cm] 

\noindent\textbf{\textsc{Semana 3}} Linux.
Introducción a UNIX: filosofía, comandos básicos. 
Editores de texto, control de procesos, redirección.
\textbf{Videos:} 1 y 2 Introducci\'on a UNIX.
Entrega ejercicio 1.
\\[-0.3cm] 

\noindent\textbf{\textsc{Semana 4}}  Linux. 

\\[-0.3cm]  

\noindent\textbf{\textsc{Semana 5}} Introducción a Python: filosofía,
sintaxis básica, operaciones aritméticas, condicionales.

\textbf{Videos:} 4 y 5 Introducci\'on a Python.

Entrega ejercicio 2.
\\[-0.3cm]   

\noindent\textbf{\textsc{Semana 6}} Operaciones con cadenas de
caracteres, listas y estructuras iterativas.

\textbf{Video:} 6  Python: Listas y strings.

Entrega ejercicio 3.
\\[-0.3cm] 

\noindent\textbf{\textsc{Semana 7}} Definición de funciones,
tipos de variables, recursividad. 

\textbf{Video:} 7 Python: Funciones, tipos de variables y recursividad.

Entrega ejercicio 4.
\\[-0.3cm] 

\noindent\textbf{\textsc{Semana 8}} Programaci\'on orientada a objetos en Python. 

\textbf{Videos:} 8 Programación Orientada a Objetos en Python.

Entrega ejercicio 5.
\\[-0.3cm]  

\noindent\textbf{\textsc{Semana 9}} Instalación (\verb+pip+) e importación de módulos. Cuadernos de iPython.

\textbf{Videos:} 9 Python: Módulos, PIP y cuadernos de Jupyter.

Entrega ejercicio 6.
\\[-0.3cm]  


\noindent\textbf{\textsc{Semana 10}} Introducción a NumPy. Arrays de numpy y operaciones entre arrays. 

\textbf{Video:} 10 Numpy con IPython

Entrega ejercicio 7.
\\[-0.3cm]  


\noindent\textbf{\textsc{Semana 11}}  Importación de datos. Visualización de datos con
\verb+matplotlib+: \verb+plot, scatter, imshow+ y  \verb+ subplot+.

\textbf{Video:} 11 import Matplotlib as plt.

Entrega ejercicio 8.
\\[-0.3cm]  

\noindent\textbf{\textsc{Semana 12}} Análisis numérico: métodos de bisección, método de Newton-Raphson. 

\textbf{Video:} 12 Métodos numéricos: ceros de funciones con IPython.

Entrega ejercicio 9.
\\[-0.3cm]  


\noindent\textbf{\textsc{Semana 13}} Introducción a
SciPy. Ajustes polinomiales y no polinomiales.

\textbf{Video:} 13 Ajuste a funciones lineales y no lineales.
Entrega ejercicio 10.
\\[-0.3cm]  

\noindent\textbf{\textsc{Semana 14}} Herramientas
estadísticas: funciones estadísticas, histogramas y ejemplos de
distribuciones.

\textbf{Video:} 14 Herramientas estadísticas con Python.

Entrega ejercicio 11.
\\[-0.3cm] 

\noindent\textbf{\textsc{Semana 15}} Métodos de Monte Carlo.


\textbf{Video:} 15 M\'etodos Monte Carlo 

Entrega ejercicio 12.
\\[-0.3cm] 

\noindent\textbf{\textsc{Semana 16}} Simulaci\'on de sistemas f\'isicos. \\[0.1cm] 


\stepcounter{mysection}
\noindent\textbf{\large \Roman{mysection} \quad Metodología}\\[-0.2cm]

\noindent\normalsize Se har\'a \'enfasis en el trabajo individual por
fuera del horarior de clase. El profesor entregar\'a una lista de
recursos que cada estudiante debe preparar {\bf antes de llegar a la
  clase presencial}.   
En este curso los temas a preparar se encuentran en videos.

Al inicio de cada clase el profesor responderá preguntas y 
hará aclaraciones sobre el material de la semana. Luego los
estudiantes resolverán individualmente los ejercicios de la semana.
Cada ejercicio semanal tiene dos partes. La primera se publica al menos un
d\'ia antes de la clase presencial y se espera que los estudiantes lo resuelvan antes de llegar a clase. La segunda se publica al
momento de la clase.
Todo el trabajo es individual y las calificaciones se otorgar\'an por
los resultados entregados exclusivamente en horario de clase por SICUA. 
Bajo ninguna circunstancia se recibirán ejercicios entregados por fuera 
del tiempo estipulado (hora oficial del final de la clase más 10 minutos).

Se recomienda traer a clase una memoria USB si le interesa guardar su 
trabajo. Se espera del estudiante que evite activamente que
su trabajo sea plagiado: cualquier sospecha de copia hará que el caso sea
llevado automáticamente a proceso disciplinario
La memoria USB también podrá ser útil si por algún motivo su computador
pierde la conexión a internet y necesita cambiar de ordenador para subir
su trabajo a la plataforma de SICUA.
\\[0.1cm] 
\newpage

\stepcounter{mysection}
\noindent\textbf{\large \Roman{mysection} \quad Calificaci\'on del curso}\\[-0.2cm]

\begin{itemize}
        \item Asistencia a clase: $10\%$
	\item 6 entregas en clase: $60\%$ ($10\%$ cada una).
	\item Examen final $30\%$.
\end{itemize}
\noindent\normalsize Durante el semestre se deben entregar 
las soluciones a los 14 ejercicios tal como indica el programa. 
Estos ejercicios se deben resolver y subir a sicuaplus antes de la
finalizaci\'on de la clase. 
Entre esos doce ejercicios se 
eligirán seis a discreci\'on del profesor que serán calificados, cada uno con un peso del $10\%$ de la nota final. 
\\[0.1cm] 

\stepcounter{mysection}
\noindent\textbf{\large \Roman{mysection} \quad Bibliografía}\\[-0.2cm]

\noindent\normalsize Bibliografía principal:

%\href{http://www.wikibooks.org}{\color{blue}\looparrowright}

\begin{itemize}
	\item H. P. Langtangen. \textit{A Primer on Scientific Programming with Python}, 2009.\\
	\href{http://link.springer.com.ezproxy.uniandes.edu.co:8080/book/10.1007\%2F978-3-642-18366-9}{\nolinkurl{http://link.springer.com.ezproxy.uniandes.edu.co:8080/book/10.1007\%2F978-3-642-18366-9}}\\[-0.6cm]
\end{itemize}

\noindent\normalsize Bibliografía complementaria:

\begin{itemize}
	\item J. V. Guttag. \textit{Introduction to Computation and Programming Using Python}, 2013.
	\item K. D. Lee. \textit{Python Programming Fundamentals}, 2011. \\
	\href{http://link.springer.com.ezproxy.uniandes.edu.co:8080/book/10.1007\%2F978-1-84996-537-8}{\nolinkurl{http://link.springer.com.ezproxy.uniandes.edu.co:8080/book/10.1007\%2F978-1-84996-537-8}}\\[-0.6cm]
	\item S. van Vugt. \textit{Beginning the Linux Command Line}, 2009.\\
	\href{http://link.springer.com.ezproxy.uniandes.edu.co:8080/book/10.1007\%2F978-1-4302-1890-6}{\nolinkurl{http://link.springer.com.ezproxy.uniandes.edu.co:8080/book/10.1007\%2F978-1-4302-1890-6}}\\[-0.6cm]
	\item G. Gr\"atzer. \textit{More Math Into  \LaTeX}, 2007.\\
	\href{http://link.springer.com.ezproxy.uniandes.edu.co:8080/book/10.1007\%2F978-0-387-68852-7}{\nolinkurl{http://link.springer.com.ezproxy.uniandes.edu.co:8080/book/10.1007\%2F978-0-387-68852-7}}
\end{itemize}

\end{document}
